\rtmtight
\begin{tblr}{width=\textwidth,colspec={|Q[c,m,2.1cm]|X[l,m]|},hlines,vlines,hspan=minimal,rowsep=1.5pt,colsep=3pt,row{1}={bg=headgray,font=\bfseries}}
\SetCell{c,m}\hfil\logo\hfil & \textbf{POLITEKNIK NEGERI MALANG}\par\textbf{JURUSAN TEKNOLOGI INFORMASI}\par\textbf{PROGRAM STUDI: D4 TEKNIK INFORMATIKA} \\
\SetCell[c=2]{l} \textbf{RENCANA TUGAS MAHASISWA (RTM) DAN RUBRIK PENILAIAN} & \\
Nama Mata Kuliah & Pemrograman Berbasis Objek \\
Kode Mata Kuliah & RTI253007 \quad \textbf{sks / Jam perminggu:} 2 SKS/4 jam \quad \textbf{Semester:} 3 \\
Dosen Pengampu & Tim Pengajar Program Studi D4 TI \\
\end{tblr}
\assessmentblock{Tugas Prinsip OOP dan Rancangan Kelas}{Case Method dan praktik pemrograman}{SCPMK0210-02501, SCPMK0704-02502}{Mahasiswa menyelesaikan tugas merancang kelas dan mengimplementasikan prinsip OOP menggunakan Java.}{Menganalisis kebutuhan, merancang diagram kelas, mengimplementasikan kode Java, menguji hasil, dan mendokumentasikan.}{Diagram kelas (UML), kode sumber Java, laporan, dan/atau bahan presentasi.}{Kelengkapan dan ketepatan rancangan kelas & 6 & Kelas, atribut, metode, relasi, dan visibilitas dirancang dengan benar sesuai prinsip OOP dan kebutuhan kasus. & Rancangan kelas sebagian besar benar, tetapi beberapa relasi atau visibilitas belum tepat. & Rancangan kelas tidak mencerminkan prinsip OOP atau tidak sesuai kebutuhan kasus. \\ \\
Implementasi enkapsulasi dan inheritance & 5 & Enkapsulasi diterapkan dengan modifier yang tepat; inheritance diimplementasikan dengan benar dan memanfaatkan reuse kode secara optimal. & Enkapsulasi atau inheritance ada tetapi penerapannya belum konsisten atau optimal. & Enkapsulasi atau inheritance tidak diterapkan atau diimplementasikan secara keliru. \\ \\
Kualitas kode dan dokumentasi & 4 & Kode bersih, terdokumentasi, nama variabel/metode deskriptif, dan mengikuti konvensi Java. & Kode berfungsi tetapi dokumentasi atau konvensi penamaan masih belum konsisten. & Kode sulit dibaca, tidak terdokumentasi, atau tidak mengikuti konvensi Java. \\}

\assessmentblock{Kuis Konsep OOP}{Kuis}{SCPMK0210-02501, SCPMK0704-02502}{Kuis tertulis untuk mengukur pemahaman konsep dasar OOP dan implementasinya menggunakan Java.}{Menjawab soal pilihan ganda dan/atau uraian terkait konsep OOP, kelas, enkapsulasi, inheritance, dan overriding.}{Jawaban tertulis kuis.}{Pemahaman konsep OOP & 5 & Konsep enkapsulasi, inheritance, dan polymorfisme dijelaskan dengan benar dan contoh penerapannya tepat. & Sebagian besar konsep dipahami tetapi contoh atau penjelasan masih kurang lengkap. & Konsep OOP tidak dipahami atau contoh penerapan keliru. \\ \\
Kemampuan analisis kode Java & 5 & Kode Java dianalisis dengan tepat, output diprediksi benar, dan error diidentifikasi secara akurat. & Sebagian besar analisis kode benar tetapi beberapa prediksi output atau identifikasi error masih keliru. & Analisis kode tidak tepat atau tidak menunjukkan pemahaman sintaks Java. \\}

\assessmentblock{UTS Evaluasi OOP Dasar}{UTS berbasis pemrograman}{SCPMK0210-02501}{Evaluasi tengah semester untuk mengukur kemampuan implementasi konsep OOP dasar menggunakan Java.}{Mengimplementasikan program Java yang menerapkan kelas, enkapsulasi, inheritance, dan relasi kelas sesuai soal.}{Kode sumber Java yang berjalan dan dokumen penjelasan desain.}{Implementasi kelas dan enkapsulasi & 5 & Kelas diimplementasikan dengan atribut private, getter/setter yang benar, dan konstruktor yang sesuai. & Implementasi kelas ada tetapi penggunaan modifier atau konstruktor belum sepenuhnya tepat. & Implementasi tidak menerapkan enkapsulasi atau struktur kelas tidak benar. \\ \\
Implementasi inheritance dan relasi kelas & 5 & Inheritance diimplementasikan dengan benar menggunakan extends, super, dan override sesuai hierarki kelas yang dirancang. & Inheritance ada tetapi penggunaan super, override, atau hierarki kelas masih belum optimal. & Inheritance tidak diterapkan atau implementasinya menghasilkan error runtime. \\}

\assessmentblock{CM Implementasi OOP Lanjut}{Case Method}{SCPMK0704-02502, SCPMK0903-02503}{Mahasiswa membangun solusi berbasis OOP menggunakan polymorfisme, interface, kelas abstrak, dan Java API untuk menyelesaikan studi kasus yang diberikan.}{Menganalisis kasus, merancang solusi OOP lanjut, mengimplementasikan kode Java, menguji fungsionalitas, dan mempresentasikan.}{Kode sumber Java, diagram arsitektur OOP, laporan, dan/atau presentasi.}{Implementasi polymorfisme dan interface & 12 & Polymorfisme diimplementasikan secara konsisten melalui interface/abstract class; metode override menghasilkan perilaku yang tepat sesuai tipe objek runtime. & Polymorfisme atau interface diimplementasikan tetapi konsistensi atau ketepatan perilaku runtime masih terbatas. & Polymorfisme atau interface tidak diimplementasikan atau menghasilkan hasil yang tidak sesuai ekspektasi. \\ \\
Penggunaan Java API dan GUI & 10 & Java API digunakan secara tepat dan efisien; komponen GUI dirancang secara logis dan responsif terhadap event pengguna. & Java API dan GUI digunakan tetapi pemilihan API atau penanganan event masih belum optimal. & Java API tidak digunakan atau GUI tidak fungsional. \\ \\
Kualitas solusi dan presentasi & 8 & Solusi memenuhi semua kebutuhan kasus, kode bersih, dan presentasi menjelaskan keputusan desain secara meyakinkan. & Solusi memenuhi sebagian besar kebutuhan tetapi presentasi atau kualitas kode masih bisa ditingkatkan. & Solusi tidak memenuhi kebutuhan kasus atau presentasi tidak dapat menjelaskan keputusan desain. \\}

\assessmentblock{UAS Aplikasi Java Terintegrasi}{UAS berbasis proyek}{SCPMK0704-02502, SCPMK0903-02503}{Mahasiswa membangun aplikasi Java lengkap berbasis GUI dan database yang menerapkan seluruh prinsip OOP secara terintegrasi.}{Merancang aplikasi, mengimplementasikan OOP lengkap, mengintegrasikan GUI dan database, menguji fungsionalitas, dan mendokumentasikan.}{Aplikasi Java yang berjalan, kode sumber, dokumentasi teknis, dan laporan pengujian.}{Kelengkapan dan fungsionalitas aplikasi & 4 & Aplikasi berjalan lengkap dengan semua fitur yang diminta, GUI responsif, dan koneksi database berfungsi dengan benar. & Sebagian besar fitur berfungsi tetapi beberapa fungsi atau koneksi database masih ada kekurangan. & Aplikasi tidak berjalan atau banyak fitur yang tidak berfungsi. \\ \\
Penerapan prinsip OOP & 3 & Seluruh prinsip OOP (enkapsulasi, inheritance, polymorfisme, abstraksi) diterapkan secara konsisten dan tepat dalam arsitektur aplikasi. & Sebagian besar prinsip OOP diterapkan tetapi beberapa tidak konsisten atau kurang optimal. & Prinsip OOP tidak diterapkan atau aplikasi menggunakan pendekatan prosedural. \\ \\
Kualitas kode dan dokumentasi & 3 & Kode terstruktur, terdokumentasi lengkap, mengikuti konvensi Java, dan dilengkapi laporan pengujian yang komprehensif. & Kode berfungsi dan terdokumentasi sebagian, tetapi struktur atau laporan pengujian masih kurang lengkap. & Kode sulit dipahami, tidak terdokumentasi, atau tidak disertai laporan pengujian. \\}

\begin{tblr}{width=\textwidth,colspec={|Q[l,m,3.2cm]|X[l,m]|},hlines,vlines,hspan=minimal,row{1-3}={bg=headgray},rowsep=1.5pt,colsep=3pt}
Jadwal Pelaksanaan: & Minggu 1--17 sesuai RPS. Setiap evaluasi memiliki RTM dan rubrik multi-kriteria. \\
Lain-Lain Yang Diperlukan: & LMS, repositori/kode sumber Java, perangkat lunak sesuai mata kuliah, dan perangkat presentasi. \\
Pustaka: & Mengacu pustaka utama dan pendukung pada bagian RPS. \\
\end{tblr}
